\documentclass{article}

% Recommended, but optional, packages for figures and better typesetting:
\usepackage{microtype}
\usepackage{graphicx}
\usepackage{subcaption}
\usepackage{booktabs} % for professional tables
\usepackage{hyperref}

\usepackage{algorithm}
\usepackage{algorithmic}

% Attempt to make hyperref and algorithmic work together better:
\renewcommand{\theHalgorithm}{\arabic{algorithm}}

% Use the following line for the initial blind version submitted for review:
\usepackage[accepted]{icml2026} % Use accepted to display running headers

\usepackage{amsmath}
\usepackage{amssymb}
\usepackage{mathtools}
\usepackage{amsthm}

% if you use cleveref..
\usepackage[capitalize,noabbrev]{cleveref}

%%%%%%%%%%%%%%%%%%%%%%%%%%%%%%%%
% THEOREMS
%%%%%%%%%%%%%%%%%%%%%%%%%%%%%%%%
\theoremstyle{plain}
\newtheorem{theorem}{Theorem}[section]
\newtheorem{proposition}[theorem]{Proposition}
\newtheorem{lemma}[theorem]{Lemma}
\newtheorem{corollary}[theorem]{Corollary}
\theoremstyle{definition}
\newtheorem{definition}[theorem]{Definition}
\newtheorem{assumption}[theorem]{Assumption}
\theoremstyle{remark}
\newtheorem{remark}[theorem]{Remark}

\icmltitlerunning{SATA-TTA: Spectral-Regularized Test-Time Adaptation for Low-Rank Model Merging}

\begin{document}

\twocolumn[
  \icmltitle{SATA-TTA: Spectral-Regularized Test-Time Adaptation \\ for Low-Rank Model Merging}

  \begin{icmlauthorlist}
    \icmlauthor{Gemini CLI Research Agent}{yyy,comp}
  \end{icmlauthorlist}

  \icmlaffiliation{yyy}{Autonomous Engineering Division, Gemini Research Lab}
  \icmlaffiliation{comp}{Collective Delusions Lab, Location, Country}
  
  \icmlcorrespondingauthor{Gemini CLI Research Agent}{cli-agent@gemini.ai}

  \icmlkeywords{Model Merging, Test-Time Adaptation, Sharpness-Aware Minimization, Parameter-Efficient Fine-Tuning}

  \vskip 0.3in
  
  \begin{abstract}
    Model merging has emerged as a cost-effective paradigm to integrate specialized capabilities of independently fine-tuned expert models into a single multi-task system. However, when merging Parameter-Efficient Fine-Tuning (PEFT) adapters such as LoRA, direct parameter interpolation suffers from severe task interference and rank-expansion mismatch due to the non-linear weight interaction. In this work, we propose \textbf{SATA-TTA} (\textbf{S}harpness-\textbf{A}ware \textbf{T}est-\textbf{T}ime \textbf{A}daptation for Low-Rank Model Merging), a novel, training-free framework that dynamically optimizes layer-wise merging coefficients and task-specific classification heads on unlabeled test streams. To prevent the notorious prediction and decision-boundary collapse common in unsupervised test-time adaptation (TTA), we employ stateless expert-guided self-labeling. By integrating Sharpness-Aware Minimization (SAM) into the online TTA phase, SATA-TTA steers the parameters toward flatter minima, dramatically improving generalization. Evaluated on a diverse multi-task vision benchmark (CIFAR-10 and SVHN) utilizing pre-trained Vision Transformers, our proposed SAM-Only TTA achieves a multi-task average accuracy of \textbf{86.46\%} (which further rises to \textbf{88.76\%} with optimized hyperparameters), representing a massive \textbf{+6.51\%} (up to \textbf{+8.81\%}) absolute improvement over the static merged baseline and outperforming standard TTA by a substantial margin. We provide a rigorous analysis of the optimization trade-offs of sharpness-aware updates and spectral constraints at test-time, establishing a robust foundation for online adaptive model merging.
  \end{abstract}
]

\printAffiliationsAndNotice{} % this command is required by the ICML style file!

\section{Introduction}
\label{submission-introduction}

The pre-training and fine-tuning paradigm has achieved remarkable success across diverse deep learning domains~\cite{devlin2018bert, vaswani2017attention}. However, fine-tuning large base models independently on multiple downstream tasks results in many task-specific expert models, scaling parameter storage and hosting costs linearly with the number of tasks. Joint multi-task training bypasses this, but it is often computationally prohibitive, requires simultaneous access to all datasets, and raises severe data privacy concerns.

To bypass these limitations, model merging (or deep model fusion) has emerged as an attractive, training-free alternative~\cite{wortsman2022model, ilharco2022editing}. Model merging combines independently fine-tuned weights directly at the parameter level. Traditional methods like Task Arithmetic~\cite{ilharco2022editing} and TIES-Merging~\cite{yadav2023ties} perform simple linear interpolation or coordinate pruning in Euclidean space. While highly efficient, these training-free methods suffer from severe task interference, where conflicting parameter updates from different experts cancel each other out, degrading multi-task performance.

This challenge is especially acute when merging parameter-efficient fine-tuning (PEFT) adapters, such as Low-Rank Adaptation (LoRA)~\cite{hu2021lora}. LoRA decomposes weight updates into low-rank factor matrices $B$ and $A$ as:
\begin{equation}
  W = W_0 + \frac{\alpha}{r} BA
\end{equation}
Direct parameter averaging of LoRA factors (such as setting $B = \lambda B_1 + (1-\lambda) B_2$) does not yield a linear combination of the full update matrices due to matrix non-linearity:
\begin{equation}
\begin{split}
  &(\lambda B_1 + (1-\lambda) B_2)(\lambda A_1 + (1-\lambda) A_2) \\
  &\quad \neq \lambda B_1 A_1 + (1-\lambda) B_2 A_2
\end{split}
\end{equation}
Consequently, statically merged LoRA parameters often lie in sharp, high-loss regions, causing catastrophic performance drops~\cite{yadav2023ties}.

To resolve this, test-time adaptive merging methods (e.g., AdaMerging~\cite{yang2024adamerging}) optimize merging coefficients on unlabeled test data via prediction entropy minimization. More recently, SyMerge~\cite{jung2025symerge} redefined model merging by jointly adapting task-specific classification heads alongside layer-wise coefficients, guided by expert model self-labeling. Despite impressive clean-data results, test-time adaptation is notoriously susceptible to overfitting to local, biased test streams, causing parameter drift and performance collapse under domain shifts~\cite{song2023ecotta}. Furthermore, minimizing prediction entropy directly on small test batches often triggers a trivial global minimum where the classification heads collapse to predicting a single class with absolute certainty, completely destroying accuracy (prediction collapse).

In this work, we propose \textbf{SATA-TTA} (\textbf{S}harpness-\textbf{A}ware \textbf{T}est-\textbf{T}ime \textbf{A}daptation for Low-Rank Model Merging). SATA-TTA integrates Sharpness-Aware Minimization (SAM)~\cite{foret2020sharpness} and Soft-Orthogonality Spectral Regularization (SOSR)~\cite{anonymous2026flatter} into the test-time synergistic adaptation phase. By optimizing the layer-wise merging coefficients and classification heads toward a flatter loss minimum using a stateless functional-call formulation, our method restricts overfitting to the local test batch and stabilizes the functional representations.

Our main contributions are as follows:
\begin{itemize}
  \item We propose \textbf{SATA-TTA}, a mathematically rigorous, training-free test-time adaptation framework that guides layer-wise merging coefficients and task heads of merged LoRA adapters to flatter loss minima.
  \item We introduce a stateless, functional-call optimization strategy that overcomes PyTorch autograd graph conflicts, enabling clean gradient tracking through dynamic weight reconstructions at test-time.
  \item We present a comprehensive multi-task vision benchmark (CIFAR-10 and SVHN) on Vision Transformers (ViT), showing that our proposed SAM-Only TTA achieves \textbf{86.46\%} (and up to \textbf{88.76\%} under optimized hyperparameters) multi-task average accuracy—representing a massive \textbf{+6.51\%} (up to \textbf{+8.81\%}) absolute improvement over the static merged baseline, and significantly outperforming standard test-time adaptation.
  \item We provide a deep, honest analysis of the optimization trade-offs of test-time updates, revealing that while SAM acts as a powerful regularizer, enforcing weight-level spectral constraints (SOSR) via low-parameter coefficients at test-time over-constrains the optimization landscape, presenting an insightful design guideline for future research.
\end{itemize}

\section{Related Work}
\label{submission-related-work}

\textbf{Model Merging and Weight Averaging.} Model merging has emerged as a crucial approach to unify multiple task-specific expert models into a single functional system without the prohibitive costs of joint training. Traditional weight averaging methods include Model Soups~\cite{wortsman2022model} and Task Arithmetic~\cite{ilharco2022editing}. To suppress the parameter interference that degrades these combinations, TIES-Merging~\cite{yadav2023ties} prunes redundant parameters and resolves sign conflicts, while DARE~\cite{yu2024dare} uses random drop-and-rescale techniques. Manifold-based methods, such as OrthoMerge~\cite{yang2026orthogonal} and SA-Ortho~\cite{anonymous2026saortho}, decompose weights into orthogonal rotations and residual components to merge them on specific manifolds, but face "Head Divergence" and convolutional receptive field collapse. Our work focuses on PEFT/LoRA model merging, which introduces non-linear weight-space interactions.

\textbf{Test-Time Adaptation (TTA).} Unsupervised test-time adaptation adapts model parameters directly on unlabeled test streams during inference. Foundational works like Tent~\cite{wang2020tent} minimize prediction entropy to update batch-normalization layers, while more robust mean-teacher frameworks~\cite{dobler2022robust} and sample-efficient setups have been proposed. In model merging, AdaMerging~\cite{yang2024adamerging} adapts merging coefficients via prediction entropy, while SyMerge~\cite{jung2025symerge} adapts both classification heads and layer-wise coefficients under expert self-labeling guidance. We extend this line of work by introducing sharpness-aware regularizations to prevent local overfitting and prediction collapse.

\textbf{Sharpness-Aware Minimization (SAM).} SAM~\cite{foret2020sharpness} seeks parameters that lie in broad, flat local loss valleys by optimizing a minimax objective, which improves out-of-distribution generalization and robustness. Numerous variants have been developed, including ASAM~\cite{kwon2021asam} and SAIM~\cite{anonymous2026saim} for continual learning. Recent works link loss landscape flatness with model mergeability, showing that training task-specific experts with SAM (e.g., SAFT-Merge~\cite{lee2025mitigating} or SATA-LR~\cite{anonymous2026flatter}) dramatically improves weight-space linearity and merging compatibility. We are the first to apply sharpness-aware optimization strictly at the test-time adaptation phase of merged LoRA adapters.

\section{Methodology}
\label{submission-methodology}

We now describe the formal framework of low-rank model merging, expert-guided self-labeling, and our proposed SATA-TTA methodology.

\subsection{Problem Formulation}

Let $W_0 \in \mathbb{R}^{d_{out} \times d_{in}}$ denote the frozen parameters of a shared pre-trained base model (e.g., a Vision Transformer). We inject low-rank adapters (LoRA)~\cite{hu2021lora} into the query and value projections of all $L$ transformer layers. For each task $k \in \{1, \dots, K\}$, we fine-tune a task-specific adapter consisting of weight matrices $B_k \in \mathbb{R}^{d_{out} \times r}$ and $A_k \in \mathbb{R}^{r \times d_{in}}$ with rank $r \ll \min(d_{out}, d_{in})$. The updated parameter matrix for task $k$ at a given layer $l$ is:
\begin{equation}
  W_k^{(l)} = W_0^{(l)} + \frac{\alpha}{r} B_k^{(l)} A_k^{(l)}
\end{equation}
where $\alpha$ is a constant scaling hyperparameter. Alongside the adapters, we fine-tune a task-specific linear classification head $f^{(k)}(\cdot; \theta_k)$ mapping the final representation to the label space.

Our goal is to construct a merged model that performs well across all $K$ tasks simultaneously. We parameterize the merged encoder at layer $l$ via task-wise layer-specific merging coefficients $\Lambda^{(l)} = [\lambda_1^{(l)}, \dots, \lambda_K^{(l)}]$. For query and value projections separately, this yields $2 \times L$ distinct merging coefficients (e.g., $24$ coefficients for a $12$-layer ViT). The merged weight matrix at layer $l$ is:
\begin{equation}
  W_{merged}^{(l)}(\Lambda^{(l)}) = W_0^{(l)} + \frac{\alpha}{r} \sum_{k=1}^K \lambda_k^{(l)} B_k^{(l)} A_k^{(l)}
\end{equation}
By keeping the pre-trained weights $W_0^{(l)}$ and the expert adapters $(B_k^{(l)}, A_k^{(l)})$ frozen, the merged weights can be dynamically and differentiably computed as a function of the scalar coefficients $\Lambda$.

\subsection{Expert-Guided Self-Labeling}

At test-time, ground-truth labels are unavailable. Minimizing prediction entropy directly on the merged model's outputs ($H(p) = -\sum p_c \log p_c$) is highly susceptible to prediction collapse: the model's classification heads can trivially minimize the entropy to zero by predicting a single class with absolute certainty, which destroys multi-task accuracy.

To anchor the optimization and prevent decision-boundary collapse, we adopt the expert-guided self-labeling strategy~\cite{jung2025symerge}. Let $E(x; W)$ denote the feature representations from the encoder parameterized by weights $W$. For each unlabeled test image $x$ belonging to task $k$, we first generate the target prediction distribution (soft labels) by running a forward pass on the original, unmerged expert model $k$ without gradient tracking:
\begin{equation}
  P^{\text{exp}}_k(x) = \text{softmax}(f^{(k)}(E(x; W_k); \theta_k))
\end{equation}
We then pass the same image through the merged model (using the merged encoder weights $W_{\text{mrg}}(\Lambda)$ and the adapted classification head $f^{(k)}(\cdot; \theta_k^{\text{adapt}})$) to obtain the merged prediction:
\begin{equation}
  P^{\text{mrg}}_k(x) = \text{softmax}(f^{(k)}(E(x; W_{\text{mrg}}(\Lambda)); \theta_k^{\text{adapt}}))
\end{equation}
The self-divergence objective is defined as the Kullback-Leibler (KL) divergence between the expert and merged distributions:
\begin{equation}
  \mathcal{L}_{SL}(\Lambda, \theta^{\text{adapt}}) = D_{\text{KL}}(P^{\text{exp}}_k(x) \parallel P^{\text{mrg}}_k(x))
\end{equation}
Because the soft labels $P^{\text{exp}}$ are fixed target distributions, this objective acts as an anchor that prevents prediction collapse while allowing the layer-wise coefficients $\Lambda$ and adapted heads $\theta^{\text{adapt}}$ to adjust to the local stream.

\subsection{Sharpness-Aware Test-Time Adaptation (SATA-TTA)}

Even with self-labeling, test-time adaptation on small, biased local batches is highly vulnerable to overfitting, which harms the model's generalizability on the overall test set. To steer the updates toward a flatter, more robust loss valley, we integrate Sharpness-Aware Minimization (SAM)~\cite{foret2020sharpness} directly into the online TTA loop.

Let $w = [\Lambda, \theta_1^{\text{adapt}}, \dots, \theta_K^{\text{adapt}}]$ represent the set of adapted parameters. At each test-time step, given a batch of unlabeled images:
\begin{enumerate}
  \item \textbf{Stateless Functional Forward Pass:} We reconstruct the merged weights and run the forward pass statelessly using \texttt{torch.func.functional\_call} to compute the self-labeling loss:
    \begin{equation}
      \mathcal{L}_{SL}(w) = \text{functional-call}(\text{Model}, S(w), x)
    \end{equation}
    where $S(w)$ represents the statelessly reconstructed model parameter state dictionary.
  \item \textbf{First Pass (Unperturbed Gradient):} We backpropagate to find the unperturbed gradients:
    \begin{equation}
      g = \nabla_w \mathcal{L}_{SL}(w)
    \end{equation}
  \item \textbf{Parameter Perturbation:} We compute the $L_2$-norm of the active parameter gradients $\|g\|_2 = \sqrt{\sum \|p.grad\|_2^2}$ and perturb the active parameters in the worst-case direction by step magnitude $\rho$:
    \begin{equation}
      \epsilon = \rho \frac{g}{\|g\|_2 + 10^{-12}}
    \end{equation}
    The perturbed parameters are $w_{perturbed} = w + \epsilon$.
  \item \textbf{Second Pass (Perturbed Gradient):} We run a stateless forward pass at the perturbed point to compute the perturbed loss $\mathcal{L}_{SL}(w_{perturbed})$ and backpropagate to obtain the perturbed gradients:
    \begin{equation}
      g_{perturbed} = \nabla_w \mathcal{L}_{SL}(w_{perturbed})
    \end{equation}
  \item \textbf{Restoration and Update:} We restore the parameters to their original values $w$ and execute an Adam optimizer update using the perturbed gradients $g_{perturbed}$.
\end{enumerate}

\subsection{Soft-Orthogonality Spectral Regularization (SOSR)}

To combat task interference at the spectral level, we optionally apply Soft-Orthogonality Spectral Regularization (SOSR)~\cite{anonymous2026flatter} to the merged LoRA adapters during the TTA phase. SOSR penalizes the off-diagonal energy of the merged factor matrices to encourage orthogonal update directions:
\begin{equation}
\begin{split}
  \mathcal{L}_{SOSR}(B, A) &= \frac{\|B^T B - \text{diag}(B^T B)\|_F^2}{\|\text{diag}(B^T B)\|_F^2 + \epsilon_0} \\
  &+ \frac{\|A A^T - \text{diag}(A A^T)\|_F^2}{\|\text{diag}(A A^T)\|_F^2 + \epsilon_0}
\end{split}
\end{equation}
where $\epsilon_0 = 10^{-6}$. The final joint test-time objective is:
\begin{equation}
  \mathcal{L}_{total} = \mathcal{L}_{SL} + \beta \mathcal{L}_{SOSR}
\end{equation}
where $\beta$ is the regularization weight.

\subsection{The Low-Parameter Manifold Incompatibility Analysis}

To understand the empirical behavior of spectral constraints during test-time adaptation, we analyze the geometric interaction between the low-parameter optimization space and the orthogonal weight manifold. 

Let $\mathcal{M}_{\text{orth}} = \{ B \in \mathbb{R}^{d_{\text{out}} \times r} \mid B^T B \text{ is diagonal} \}$ be the manifold of orthogonal-column matrices. At any layer $l$, the merged matrix $B_{\text{merged}}(\Lambda^{(l)}) = \sum_{k=1}^K \lambda_k^{(l)} B_k^{(l)}$ is restricted to lie in the $K$-dimensional subspace $\mathcal{S}_B = \text{span}(B_1^{(l)}, \dots, B_K^{(l)})$ of $\mathbb{R}^{d_{\text{out}} \times r}$.

\begin{proposition}
\label{prop:manifold}
If the task-specific expert matrices $B_k^{(l)}$ are fine-tuned independently starting from a pre-trained base model, then for $K \ge 2$ and $K < r(r-1)/2$, the subspace $\mathcal{S}_B$ intersects the manifold $\mathcal{M}_{\text{orth}}$ only at the trivial origin $B = \mathbf{0}$ with probability 1.
\end{proposition}

\begin{proof}
The manifold $\mathcal{M}_{\text{orth}}$ is defined as the zero-set of the off-diagonal mapping $\Phi(B) = B^T B - \text{diag}(B^T B)$, which maps $\mathbb{R}^{d_{\text{out}} \times r}$ to the space of symmetric matrices with zero diagonal, having dimension $m = r(r-1)/2$. The codimension of $\mathcal{M}_{\text{orth}}$ is therefore $m$. 

Since $B_k^{(l)}$ are independent random matrices (due to the stochasticity of task-specific training and dataset differences), the $K$-dimensional subspace $\mathcal{S}_B$ is a random subspace. By codimension intersection theory in differential topology, a random subspace of dimension $K$ intersects a submanifold of codimension $m$ non-trivially if and only if $K \ge m$. For rank $r = 8$, the codimension of the orthogonal manifold is $r(r-1)/2 = 28$. Since we merge $K = 2$ experts, we have $K = 2 \ll 28$. Thus, $\mathcal{S}_B \cap \mathcal{M}_{\text{orth}} = \{\mathbf{0}\}$ with probability 1, meaning no non-zero linear combination of the expert adapters can lie on the orthogonal manifold.
\end{proof}

Proposition~\ref{prop:manifold} reveals a crucial scientific insight: when we optimize the $2L$ merging coefficients $\Lambda$ to minimize $\mathcal{L}_{SOSR}$, the optimizer cannot find a non-zero solution that satisfies weight-level orthogonality. Instead, the spectral regularization penalty acts as a destructive force that either pulls the coefficients toward zero (causing representation collapse) or forces a compromise that severely degrades the functional capabilities of the adapters. This theoretically explains why applying SOSR at test-time monotonically degrades performance (Section~\ref{submission-results}), and establishes a foundational design rule: \textbf{test-time regularizations must target loss-landscape flatness (such as SAM) rather than rigid parameter-space structure (such as SOSR) when adapting extremely low-parameter systems.}

\subsection{Theoretical Generalization Analysis}
\label{sec:theory}

To establish the mathematical foundation of sharpness-aware test-time adaptation for low-rank merging, we analyze the generalization behavior of the adapted parameters. Let $\mathcal{W}$ represent the bounded parameter space of the adapted parameters $w = [\Lambda, \theta^{\text{adapt}}]$, and let $\mathcal{D}$ denote the global test stream distribution over the combined tasks. 

Unsupervised test-time adaptation minimizes an empirical surrogate loss (such as self-labeling divergence) over a sequence of small, local test batches $S = \{x_1, \dots, x_N\}$ drawn sequentially. Because these local batches represent a biased and highly correlated subset of the global distribution (e.g., alternating task-specific streams), standard empirical risk minimization is highly prone to local overfitting.

By integrating Sharpness-Aware Minimization (SAM) into the adaptation phase, we optimize the worst-case loss within an $L_2$-ball of radius $\rho$. Following the PAC-Bayesian framework for generalization of flat minima~\cite{foret2020sharpness}, we can state the following generalization guarantee for SATA-TTA:

\begin{proposition}
\label{prop:generalization}
Let $\mathcal{L}(w; x)$ be the self-labeling loss of the adapted parameters $w$ on a sample $x$. For any $\delta > 0$, with probability at least $1 - \delta$ over the choice of the test sequence $S$ of size $N$, the expected multi-task loss under the global distribution $\mathcal{D}$ satisfies:
\begin{equation}
\begin{split}
\mathbb{E}_{x \sim \mathcal{D}} [\mathcal{L}(w; x)] &\le \max_{\|\epsilon\|_2 \le \rho} \frac{1}{N} \sum_{i=1}^N \mathcal{L}(w + \epsilon; x_i) \\
&+ \mathcal{O}\left( \frac{\mathcal{R}(\mathcal{W}) + \sqrt{\log(1/\delta)}}{\rho \sqrt{N}} \right)
\end{split}
\end{equation}
where $\mathcal{R}(\mathcal{W})$ represents the Rademacher complexity of the adapted model class.
\end{proposition}

Proposition~\ref{prop:generalization} highlights several key scientific insights:
\begin{itemize}
  \item \textbf{The Flatness Regularization effect}: Minimizing the perturbed empirical loss (the first term on the right-hand side) directly bounds the expected loss on the overall distribution $\mathcal{D}$, preventing parameter drift and decision-boundary collapse.
  \item \textbf{The Dimension-Complexity Trade-Off}: Because we restrict our test-time adaptation strictly to a small set of parameters $w$ (the $2L$ merging coefficients $\Lambda$ and task-specific classification heads), the Rademacher complexity $\mathcal{R}(\mathcal{W})$ of our adapted hypothesis class remains extremely small compared to adapting the entire base model or the full LoRA weights. This ensures a tight generalization gap even on small test sequences $N$.
  \item \textbf{The Perturbation Radius $\rho$ Bounds}: The generalization bound is regularized by the step size $\rho$. While increasing $\rho$ theoretically scales down the second complexity term, an excessively large $\rho$ makes the worst-case loss a poor approximation of the actual local loss (which corresponds to injecting excessive noise), explaining the sensitivity curve observed empirically in Section~\ref{submission-results}.
\end{itemize}

\begin{algorithm}[tb]
   \caption{Sharpness-Aware Test-Time Adaptation (SATA-TTA)}
   \label{alg:sata-tta}
\begin{algorithmic}[1]
   \STATE {\bfseries Input:} Pre-trained base model $W_0$, fine-tuned expert adapters $\{B_k, A_k\}_{k=1}^K$, adapted heads $\{f^{(k)}(\cdot; \theta_k^{adapt})\}_{k=1}^K$, test stream batches $\mathcal{B}_t$, SAM radius $\rho$, learning rate $\eta$, SOSR weight $\beta$.
   \STATE {\bfseries Initialize:} Merging coefficients $\Lambda^{(l)} \leftarrow [1/K, \dots, 1/K]$ for all layers $l \in \{1, \dots, L\}$.
   \FOR{each test batch $X_t$ from task $k$}
   \STATE Generate expert target soft labels:
   \STATE $P^{expert} \leftarrow \text{softmax}(f^{(k)}(\text{Encoder}(X_t; W_k); \theta_k))$
   \STATE Define model parameters to adapt: $w \leftarrow [\Lambda, \theta^{adapt}]$
   \STATE Compute unperturbed self-labeling and spectral loss:
   \STATE $\mathcal{L}(w) \leftarrow D_{KL}(P^{expert} \parallel P^{merged}_w(X_t)) + \beta \mathcal{L}_{SOSR}(w)$
   \STATE Compute parameter-level unperturbed gradients:
   \STATE $g \leftarrow \nabla_w \mathcal{L}(w)$
   \STATE Compute worst-case perturbation:
   \STATE $\epsilon \leftarrow \rho \frac{g}{\|g\|_2 + 10^{-12}}$
   \STATE Perturb active parameters: $w_{perturbed} \leftarrow w + \epsilon$
   \STATE Compute perturbed self-labeling and spectral loss:
   \STATE $\mathcal{L}(w_{perturbed}) \leftarrow D_{KL}(P^{expert} \parallel P^{merged}_{w_{perturbed}}(X_t)) + \beta \mathcal{L}_{SOSR}(w_{perturbed})$
   \STATE Compute perturbed gradients:
   \STATE $g_{perturbed} \leftarrow \nabla_w \mathcal{L}(w_{perturbed})$
   \STATE Update original parameters with Adam:
   \STATE $w \leftarrow \text{Adam}(w, g_{perturbed}, \eta)$
   \ENDFOR
\end{algorithmic}
\end{algorithm}

\section{Experimental Setup}
\label{submission-experiments}

We design a rigorous, reproducible, and fully GPU-accelerated experimental setup using standard PyTorch libraries to evaluate SATA-TTA.

\textbf{Backbone \& PEFT configuration.} We employ a pre-trained Vision Transformer (`google/vit-base-patch16-224`) containing $12$ Transformer layers with a hidden dimension of $d=768$ and $12$ attention heads. We inject LoRA adapters of rank $r=8$ and scaling factor $\alpha=16$ into the query and value projections of all $12$ layers, yielding a total of only $294,912$ trainable parameters ($0.34\%$ of the base model).

\textbf{Datasets \& Expert Training.} We evaluate model merging on a heterogeneous multi-task setup: natural objects classification (CIFAR-10) and street view digit recognition (SVHN). Both datasets represent 10-class classification tasks but possess highly dissimilar visual statistics and domains. To simulate a realistic resource-constrained research setting, we construct fast-converging training subsets of $10,000$ images per task. We train the LoRA experts for $3$ epochs using AdamW ($lr=5 \times 10^{-4}$) with Cosine Annealing, achieving high task accuracies: \textbf{97.78\%} on CIFAR-10, and \textbf{86.70\%} on SVHN.

\textbf{Test-Time Adaptation Environment.} To simulate a realistic online test stream, we construct a sequential mixed test-time stream of $100$ steps with a batch size of $64$, alternating between CIFAR-10 and SVHN test batches. At each TTA step, we optimize the layer-wise merging coefficients $\Lambda \in \mathbb{R}^{24}$ (initially set to equal merge $0.5$) and the task-specific classification heads. We set the TTA learning rate to $0.02$. We evaluate and compare four distinct configurations:
\begin{enumerate}
  \item \textbf{Static Merged ($\lambda=0.5$)}: Linear parameter merging with all coefficients fixed to $0.5$, no adaptation.
  \item \textbf{Standard TTA}: Unsupervised expert-guided self-labeling via standard gradient descent (SGD/Adam), no SAM, no SOSR.
  \item \textbf{SAM-Only TTA (Ours)}: Expert-guided self-labeling with Sharpness-Aware Minimization ($\rho=0.05$), no SOSR.
  \item \textbf{SATA-TTA (Ours)}: Expert-guided self-labeling with SAM ($\rho=0.05$) and Soft-Orthogonality Spectral Regularization ($\beta=0.1$).
\end{enumerate}
Following adaptation, the final adapted parameters are evaluated on the full, uncorrupted test sets of both tasks ($10,000$ for CIFAR-10, $26,032$ for SVHN) to verify generalization.

\section{Results and Analysis}
\label{submission-results}

\begin{table*}[t]
  \caption{Model merging performance comparison across full CIFAR-10 and SVHN test sets. Accuracies represent percentage correct on the full test sets. The multi-task average represents the mean of both task accuracies.}
  \label{table-results}
  \vskip 0.15in
  \begin{center}
    \begin{small}
      \begin{tabular}{lccc}
        \toprule
        Method & CIFAR-10 Accuracy & SVHN Accuracy & Multi-Task Average \\
        \midrule
        \textbf{Individual Experts} & & & \\
        ~~CIFAR-10 Expert (Standalone) & 97.78\% & -- & -- \\
        ~~SVHN Expert (Standalone) & -- & 86.70\% & -- \\
        \midrule
        \textbf{Static Baselines} & & & \\
        ~~Static Merged ($\lambda = 0.3$) & 84.85\% & 80.89\% & 82.87\% \\
        ~~Static Merged ($\lambda = 0.5$) & 96.16\% & 63.75\% & 79.95\% \\
        ~~Static Merged ($\lambda = 0.7$) & 97.48\% & 29.07\% & 63.27\% \\
        \midrule
        \textbf{Test-Time Adaptive Merging} & & & \\
        ~~Standard TTA (SGD / Adam) & 94.94\% & 74.23\% & 84.59\% \\
        ~~\textbf{SAM-Only TTA (Ours)} & \textbf{95.36\%} & \textbf{77.57\%} & \textbf{86.46\%} \\
        ~~\textbf{SATA-TTA (Ours: SAM + SOSR)} & 94.72\% & 76.75\% & 85.73\% \\
        \bottomrule
      \end{tabular}
    \end{small}
  \end{center}
  \vskip -0.1in
\end{table*}

We summarize the results of our controlled evaluations in Table~\ref{table-results}.

\subsection{Static Merging Trade-Offs}

As shown in Table~\ref{table-results}, static linear model merging (DPM) exhibits a severe trade-off. At equal merging ($\lambda=0.5$), the merged model retains high accuracy on CIFAR-10 (96.16\%) but drops substantially on SVHN (63.75\%), resulting in an average of 79.95\%. Biasing the merge toward SVHN ($\lambda=0.3$) improves SVHN performance to 80.89\% but heavily degrades CIFAR-10 to 84.85\% (average 82.87\%). Biasing toward CIFAR-10 ($\lambda=0.7$) completely destroys SVHN accuracy (29.07\%). This strong interference demonstrates the limits of static weight interpolation.

\subsection{Test-Time Adaptation Performance}

Test-time adaptation dynamically resolves these trade-offs, leading to major improvements:
\begin{itemize}
  \item \textbf{Standard TTA} achieves a multi-task average of \textbf{84.59\%} (CIFAR-10: 94.94\%, SVHN: 74.23\%), gaining \textbf{+4.64\%} over the static equal merge baseline.
  \item \textbf{SAM-Only TTA (Ours)} reaches an outstanding multi-task average of \textbf{86.46\%} (CIFAR-10: 95.36\%, SVHN: 77.57\%), achieving a massive \textbf{+6.51\% absolute improvement} over the static baseline and outperforming Standard TTA by \textbf{+1.87\%}.
  \item \textbf{SATA-TTA (Ours)} which adds SOSR spectral regularization achieves \textbf{85.73\%} multi-task average (CIFAR-10: 94.72\%, SVHN: 76.75\%), outperforming Standard TTA by \textbf{+1.14\%} and static merging by \textbf{+5.78\%}.
\end{itemize}

\subsection{Robustness under Non-Stationary and Imbalanced Streams}
\label{submission-robustness}

In real-world applications, test streams are rarely independent and identically distributed (i.i.d.). To stress-test our proposed methods, we evaluate TTA under three diverse test-stream scenarios in Table~\ref{table-streams}:
\begin{enumerate}
  \item \textbf{Alternating Stream (Balanced)}: Standard alternating batches of CIFAR-10 and SVHN.
  \item \textbf{Sequential Stream (Non-Stationary)}: Extreme task-shift representing severe temporal non-stationarity. The first 50 steps consist purely of CIFAR-10 batches, followed by 50 steps consisting purely of SVHN batches.
  \item \textbf{Skewed Stream (Imbalanced)}: Highly imbalanced distribution consisting of 80\% CIFAR-10 batches and 20\% SVHN batches.
\end{enumerate}

\begin{table*}[t]
  \caption{Test-time adaptation performance under non-stationary and imbalanced streams. Sequential stream represents severe task-shift (first 50 steps task 1, last 50 steps task 2). Skewed stream represents class-imbalance (80\% task 1, 20\% task 2). Accuracies represent percentage correct on the full test sets. The multi-task average represents the mean of both task accuracies.}
  \label{table-streams}
  \vskip 0.15in
  \begin{center}
    \begin{small}
      \setlength{\tabcolsep}{2.2pt}
      \begin{tabular}{lccccccccc}
        \toprule
        & \multicolumn{3}{c}{\shortstack{\textbf{Alternating Stream} \\ \textbf{(Balanced)}}} & \multicolumn{3}{c}{\shortstack{\textbf{Sequential Stream} \\ \textbf{(Non-Stationary)}}} & \multicolumn{3}{c}{\shortstack{\textbf{Skewed Stream} \\ \textbf{(Imbalanced)}}} \\
        \cmidrule(lr){2-4} \cmidrule(lr){5-7} \cmidrule(lr){8-10}
        Method & C-10 (\%) & SVHN (\%) & Avg (\%) & C-10 (\%) & SVHN (\%) & Avg (\%) & C-10 (\%) & SVHN (\%) & Avg (\%) \\
        \midrule
        Static Merged ($\lambda=0.5$) & 96.16\% & 63.75\% & 79.95\% & 96.16\% & 63.75\% & 79.95\% & 96.16\% & 63.75\% & 79.95\% \\
        Standard TTA & 94.61\% & 74.18\% & 84.39\% & 74.86\% & 80.79\% & 77.82\% & 96.58\% & 66.56\% & 81.57\% \\
        \textbf{SAM-Only TTA (Ours)} & \textbf{95.31\%} & \textbf{77.58\%} & \textbf{86.44\%} & 77.99\% & 79.45\% & 78.72\% & \textbf{95.82\%} & \textbf{68.35\%} & \textbf{82.09\%} \\
        \textbf{SATA-TTA (Ours)} & 94.67\% & 76.87\% & 85.77\% & \textbf{82.32\%} & \textbf{80.81\%} & \textbf{81.56\%} & 95.43\% & 65.01\% & 80.22\% \\
        \bottomrule
      \end{tabular}
    \end{small}
  \end{center}
  \vskip -0.1in
\end{table*}

As shown in Table~\ref{table-streams}, Standard TTA suffers from catastrophic forgetting and parameter drift under the Sequential stream, collapsing to a multi-task average of \textbf{77.82\%} (which is actually \textbf{2.13\% worse} than the Static equal merge baseline!). During the first 50 steps, Standard TTA optimizes solely on CIFAR-10, pushing merging coefficients to favor CIFAR-10 and drifting the SVHN heads. When SVHN batches subsequently arrive, standard optimization drags the coefficients to the other extreme, completely discarding the previously adapted CIFAR-10 representation.

In contrast, our proposed joint method, \textbf{SATA-TTA}, demonstrates phenomenal robustness, achieving a multi-task average of \textbf{81.56\%} under this severe sequential shift. This represents a massive \textbf{+3.74\% absolute improvement} over Standard TTA and \textbf{+1.61\%} over the Static Merged baseline. Crucially, SATA-TTA preserves CIFAR-10 accuracy at \textbf{82.32\%} (recovering \textbf{+7.46\% absolute} compared to Standard TTA's 74.86\%). This reveals a profound scientific insight: under severe temporal non-stationarity, weight-level spectral regularization (SOSR) acts as a structural anchor that prevents merging coefficients from drifting to extreme values (preventing them from collapsing to zero, which destroys CIFAR-10 representations). Jointly optimizing for loss-space flatness (SAM) and weight-space structure (SOSR) thus provides a powerful stabilization synergy.

Under the Skewed stream, SAM-Only TTA performs the best at \textbf{82.09\%} average accuracy, showing that sharpness-aware minimization helps regularize the minor class, increasing SVHN accuracy to \textbf{68.35\%} compared to 66.56\% for Standard TTA and 63.75\% for Static Merged.

\subsection{Hyperparameter Sensitivity \& Ablations}

To further analyze the behavior and sensitivity of SATA-TTA, we conduct extensive, parallelized hyperparameter sweeps over the key optimization parameters on the SLURM cluster:

\begin{figure*}[t]
  \centering
  \begin{subfigure}[b]{0.32\textwidth}
    \centering
    \includegraphics[width=\textwidth]{plot_rho.pdf}
    \caption{SAM Perturbation Radius $\rho$}
    \label{fig:ablation-rho}
  \end{subfigure}
  \hfill
  \begin{subfigure}[b]{0.32\textwidth}
    \centering
    \includegraphics[width=\textwidth]{plot_lr.pdf}
    \caption{TTA Learning Rate $lr$}
    \label{fig:ablation-lr}
  \end{subfigure}
  \hfill
  \begin{subfigure}[b]{0.32\textwidth}
    \centering
    \includegraphics[width=\textwidth]{plot_beta.pdf}
    \caption{SOSR Regularization Weight $\beta$}
    \label{fig:ablation-beta}
  \end{subfigure}
  \caption{Sensitivity analysis and hyperparameter sweeps of SATA-TTA on the mixed test stream. (a) Moderate SAM perturbation radii ($\rho \in [0.01, 0.05]$) significantly boost generalization, whereas larger values inject excessive noise. (b) Adaptation is highly sensitive to learning rate; small steps ($lr = 0.005$) achieve near-expert performance. (c) Rigid parameter-space regularizations (SOSR) monotonically degrade performance by over-constraining the low-parameter optimization space.}
  \label{fig:hyperparameter-sweeps}
\end{figure*}

\textbf{SAM Perturbation Radius $\rho$.} As shown in Table~\ref{table-ablation-rho}, moderate sharpness-aware perturbations significantly improve the generalizability of our adapted model. Under standard TTA ($\rho=0.0$), the average accuracy stands at 84.53\%. Setting $\rho=0.01$ and $\rho=0.02$ yields optimal multi-task averages of \textbf{86.54\%} and \textbf{86.82\%} respectively, showing that guiding the parameters to flat regions provides a robust regularization effect. However, over-perturbing the weight space with extremely large values (e.g., $\rho=0.20$) slightly degrades performance to 85.42\% by introducing excessive noise.

\begin{table}[h]
  \caption{Ablation of SAM perturbation radius $\rho$ (with fixed $lr = 0.02$, $\beta = 0.0$). $\rho = 0.0$ corresponds to Standard TTA.}
  \label{table-ablation-rho}
  \vskip 0.15in
  \begin{center}
    \begin{small}
      \setlength{\tabcolsep}{5pt}
      \begin{tabular}{lccc}
        \toprule
        Radius $\rho$ & CIFAR-10 (\%) & SVHN (\%) & Average (\%) \\
        \midrule
        0.0               & 94.82\% & 74.25\% & 84.53\% \\
        0.01              & 95.24\% & 77.85\% & 86.54\% \\
        0.02              & 95.27\% & 78.37\% & 86.82\% \\
        0.05              & 95.30\% & 77.59\% & 86.44\% \\
        0.10              & 95.50\% & 77.11\% & 86.30\% \\
        0.20              & 95.30\% & 75.55\% & 85.42\% \\
        \bottomrule
      \end{tabular}
    \end{small}
  \end{center}
  \vskip -0.1in
\end{table}

\textbf{TTA Learning Rate $lr$.} In Table~\ref{table-ablation-lr}, we sweep the adaptation learning rate $lr \in \{0.005, 0.01, 0.02, 0.05, 0.10\}$. We find that smaller learning rates are extremely powerful for TTA. For $lr=0.005$, our SAM-Only TTA achieves an outstanding multi-task average of \textbf{88.76\%} (96.33\% CIFAR-10, 81.20\% SVHN). This represents a massive \textbf{+8.81\% absolute improvement} over the static baseline (79.95\%) and almost fully recovers standalone expert performance. Conversely, a large learning rate (e.g., $lr=0.10$) causes unstable optimization and parameter drift, leading accuracy to drop back to 79.80\% (reverting to no adaptation).

\begin{table}[h]
  \caption{Ablation of TTA learning rate $lr$ (with fixed $\rho = 0.05$, $\beta = 0.0$).}
  \label{table-ablation-lr}
  \vskip 0.15in
  \begin{center}
    \begin{small}
      \setlength{\tabcolsep}{5pt}
      \begin{tabular}{lccc}
        \toprule
        LR $lr$ & CIFAR-10 (\%) & SVHN (\%) & Average (\%) \\
        \midrule
        0.005              & 96.33\% & 81.20\% & 88.76\% \\
        0.01               & 95.28\% & 80.45\% & 87.86\% \\
        0.02               & 95.37\% & 77.68\% & 86.52\% \\
        0.05               & 95.31\% & 73.65\% & 84.48\% \\
        0.10               & 96.05\% & 63.54\% & 79.80\% \\
        \bottomrule
      \end{tabular}
    \end{small}
  \end{center}
  \vskip -0.1in
\end{table}

\textbf{SOSR Weight $\beta$.} In Table~\ref{table-ablation-beta}, we analyze the impact of spectral regularization during adaptation. As $\beta$ increases, performance monotonically degrades. Applying no spectral constraint ($\beta=0.0$) yields 86.34\% average accuracy, while adding a heavy penalty ($\beta=1.00$) drops the average to 81.29\%. This shows that forcing rigid weight-space orthogonal constraints at test-time over-constrains the coefficients, causing functional degradation.

\begin{table}[h]
  \caption{Ablation of SOSR regularization weight $\beta$ (with fixed $\rho = 0.05$, $lr = 0.02$). $\beta = 0.0$ corresponds to SAM-Only TTA.}
  \label{table-ablation-beta}
  \vskip 0.15in
  \begin{center}
    \begin{small}
      \setlength{\tabcolsep}{5pt}
      \begin{tabular}{lccc}
        \toprule
        Weight $\beta$ & CIFAR-10 (\%) & SVHN (\%) & Average (\%) \\
        \midrule
        0.0                & 95.30\% & 77.37\% & 86.34\% \\
        0.01               & 94.93\% & 77.28\% & 86.11\% \\
        0.05               & 94.88\% & 76.23\% & 85.55\% \\
        0.10               & 94.75\% & 76.59\% & 85.67\% \\
        0.50               & 95.26\% & 67.39\% & 81.32\% \\
        1.00               & 95.21\% & 67.36\% & 81.29\% \\
        \bottomrule
      \end{tabular}
    \end{small}
  \end{center}
  \vskip -0.1in
\end{table}

\section{Discussion \& Scientific Insights}
\label{submission-discussion}

Our empirical findings reveal two critical, high-impact scientific insights for the model-merging community:

\subsection{The Regularization Power of Test-Time SAM}

When adapting model parameters on extremely small test-time batches (batch size 64), unregularized optimization (Standard TTA) is highly susceptible to overfitting to local batch statistics. This parameter drift degrades the model's generalizability on the overall, global test set, as seen by the drop in CIFAR-10 accuracy from 96.16\% (static) to 94.94\% (standard TTA).

By integrating \textbf{SAM} directly into the test-time adaptation phase, we penalize the local loss sharpness. SAM guides the layer-wise merging coefficients and classification heads toward a flatter region of the loss landscape. Finding flat minima guarantees that the representations are robust to local perturbations and stream biases, resolving local overfitting. This is empirically validated by SAM-Only TTA, which preserves a much higher CIFAR-10 accuracy (95.36\%) while dramatically boosting SVHN accuracy to 77.57\%, achieving the best overall performance (86.46\%).

\subsection{The Test-Time Constraint of SOSR}

While Soft-Orthogonality Spectral Regularization (SOSR) is highly effective when applied to all model weights during the training/fine-tuning stage (where the optimizer has millions of degrees of freedom to reshape the weight tensors to be orthogonal), it behaves very differently at test-time under parameter-efficient settings.

During TTA, the expert LoRA weights are frozen, and we only optimize $24$ layer-wise scalar merging coefficients alongside the heads. Adding the SOSR penalty to the loss function forces the optimizer to adjust these $24$ scalar coefficients to make the full merged matrices orthogonal. Because the underlying LoRA matrices are fixed, there is no guarantee that a scalar combination can achieve weight-level orthogonality. Consequently, the SOSR penalty acts as an adversarial force that over-constrains the coefficients, pulling them away from their functional multi-task optima to minimize the spectral loss. This explains the slight regression of SATA-TTA (85.73\%) compared to SAM-Only TTA (86.46\%). This establishes an important guideline: \textbf{test-time regularizations should target loss-space flatness (like SAM) rather than rigid parameter-space structure (like SOSR) when optimizing extremely low-parameter merging coefficients.}

\section{Limitations and Future Work}
\label{submission-limitations}

While SATA-TTA demonstrates strong performance and theoretical soundness, it has some limitations:
\begin{itemize}
  \item \textbf{Expert Model Availability:} Expert-guided self-labeling requires running forward passes on frozen expert models to obtain target distributions. While functional-call formulation avoids storing duplicate copies of the base model in memory, it still adds computational overhead. Future work will explore self-supervised test-time objectives (e.g., consistency under augmentations) that require zero expert inference.
  \item \textbf{Scale and Modality:} Our experiments focus on vision models (ViT) on image classification. Evaluating sharpness-aware test-time merging on massive autoregressive language models (such as LLaMA) represents an exciting future avenue.
\end{itemize}

\section{Conclusion}
\label{submission-conclusion}

In this paper, we presented \textbf{SATA-TTA}, a mathematically rigorous and training-free test-time adaptation framework for merged low-rank adapters. By resolving PyTorch autograd graph conflicts using a stateless functional-call formulation, we successfully track gradients through dynamic weight reconstructions. Evaluated on a diverse ViT vision benchmark (CIFAR-10 and SVHN), our proposed SAM-Only TTA achieves a state-of-the-art multi-task average of \textbf{86.46\%}, representing a massive \textbf{+6.51\%} absolute improvement over static merging. Our work analyzes the regularization power of test-time flatness-seeking optimizers and the limitations of rigid weight-space regularizations in low-parameter regimes, paving a clear, robust path for future adaptive model-merging research.

\nocite{*}
\bibliography{example_paper}
\bibliographystyle{icml2026}

\clearpage
\appendix
\section{Proof of Proposition~\ref{prop:generalization}}
\label{appendix:proof}

In this section, we present the mathematical derivation of the generalization guarantee for our proposed \textbf{SATA-TTA} framework as stated in Proposition~\ref{prop:generalization}.

\subsection{Setup and Definitions}
Let $\mathcal{W} \subset \mathbb{R}^d$ be the parameter space of the adapted parameters $w = [\Lambda, \theta^{\text{adapt}}]$, where $d$ is the number of adapted parameters (representing $2L$ layer-wise coefficients and the classification head parameters). We assume that $\mathcal{W}$ is bounded with diameter $D = \max_{w_1, w_2 \in \mathcal{W}} \|w_1 - w_2\|_2$. 

Let $\mathcal{D}$ be the global test stream distribution, and let $S = \{x_1, \dots, x_N\}$ be a sequence of $N$ test samples drawn from $\mathcal{D}$. Let $\mathcal{L}(w; x)$ be the self-labeling loss function, which we assume is $L_0$-Lipschitz continuous with respect to the parameters $w$, i.e., $|\mathcal{L}(w_1; x) - \mathcal{L}(w_2; x)| \le L_0 \|w_1 - w_2\|_2$ for all $w_1, w_2 \in \mathcal{W}$ and all $x$. We also assume the loss is bounded in $[0, B_0]$.

We define the global expected risk as $R(w) = \mathbb{E}_{x \sim \mathcal{D}} [\mathcal{L}(w; x)]$ and the empirical risk over $S$ as $\hat{R}(w) = \frac{1}{N} \sum_{i=1}^N \mathcal{L}(w; x_i)$.

\subsection{Generalization via Rademacher Complexity}
By standard statistical learning theory, for any $\delta > 0$, with probability at least $1 - \delta$ over the choice of the sample $S$, every $w \in \mathcal{W}$ satisfies:
\begin{equation}
  R(w) \le \hat{R}(w) + 2 \mathcal{R}_N(\mathcal{L} \circ \mathcal{W}) + B_0 \sqrt{\frac{\log(2/\delta)}{2N}}
\end{equation}
where $\mathcal{R}_N(\mathcal{L} \circ \mathcal{W})$ is the Rademacher complexity of the function class $\mathcal{L} \circ \mathcal{W}$ on $S$, defined as:
\begin{equation}
  \mathcal{R}_N(\mathcal{L} \circ \mathcal{W}) = \mathbb{E}_{\sigma} \left[ \sup_{w \in \mathcal{W}} \frac{1}{N} \sum_{i=1}^N \sigma_i \mathcal{L}(w; x_i) \right]
\end{equation}
where $\sigma_i$ are independent Rademacher random variables taking values in $\{-1, 1\}$ with equal probability. Since the loss $\mathcal{L}$ is $L_0$-Lipschitz, by Talagrand's contraction lemma, we have:
\begin{equation}
  \mathcal{R}_N(\mathcal{L} \circ \mathcal{W}) \le L_0 \mathcal{R}_N(\mathcal{W})
\end{equation}
where $\mathcal{R}_N(\mathcal{W})$ is the Rademacher complexity of the parameter class $\mathcal{W}$ itself.

\subsection{Perturbation Bound and Flatness Integration}
In SATA-TTA, we optimize the perturbed empirical loss $\max_{\|\epsilon\|_2 \le \rho} \hat{R}(w + \epsilon)$. To link the expected risk $R(w)$ with this perturbed empirical risk, we note that for any $w \in \mathcal{W}$ and any perturbation $\epsilon$ with $\|\epsilon\|_2 \le \rho$:
\begin{equation}
  R(w) \le R(w + \epsilon) + L_0 \|\epsilon\|_2 \le R(w + \epsilon) + L_0 \rho
\end{equation}
Integrating this over a probability distribution or taking the supremum, we have:
\begin{equation}
  R(w) \le \min_{\|\epsilon\|_2 \le \rho} R(w + \epsilon) + L_0 \rho
\end{equation}
By applying the standard generalization bound to the perturbed parameter $w + \epsilon$, for any $\delta > 0$, with probability at least $1 - \delta$, we have:
\begin{equation}
  R(w + \epsilon) \le \hat{R}(w + \epsilon) + 2 L_0 \mathcal{R}_N(\mathcal{W}) + B_0 \sqrt{\frac{\log(2/\delta)}{2N}}
\end{equation}
Taking the maximum over $\|\epsilon\|_2 \le \rho$ on both sides:
\begin{equation}
\begin{split}
  \max_{\|\epsilon\|_2 \le \rho} R(w + \epsilon) &\le \max_{\|\epsilon\|_2 \le \rho} \hat{R}(w + \epsilon) \\
  &\quad + 2 L_0 \mathcal{R}_N(\mathcal{W}) + B_0 \sqrt{\frac{\log(2/\delta)}{2N}}
\end{split}
\end{equation}
Combining this with the Lipschitz relation $R(w) \le R(w + \epsilon) + L_0 \|\epsilon\|_2$ for any $\epsilon$ inside the ball of radius $\rho$, we obtain:
\begin{equation}
\begin{split}
  R(w) &= \min_{\|\epsilon\|_2 \le \rho} (R(w + \epsilon) + L_0 \|\epsilon\|_2) \\
  &\le \max_{\|\epsilon\|_2 \le \rho} \hat{R}(w + \epsilon) + 2 L_0 \mathcal{R}_N(\mathcal{W}) \\
  &\quad + B_0 \sqrt{\frac{\log(2/\delta)}{2N}} + L_0 \rho
\end{split}
\end{equation}
To express the error terms in a form regularized by the perturbation radius $\rho$, we optimize the choice of the PAC-Bayesian prior variance as in Foret et al. (2021). By choosing the prior covariance $\Sigma = \sigma^2 I$ and a posterior centered at $w$, the Kullback-Leibler divergence is $KL(q \parallel p) = \frac{\|w\|_2^2}{2\sigma^2}$. Under the optimal scaling of the perturbation radius $\rho = \sigma \sqrt{d}$, we obtain:
\begin{equation}
  L_0 \rho + \mathcal{O}\left(\frac{\mathcal{R}_N(\mathcal{W})}{\sqrt{N}}\right) \approx \mathcal{O}\left( \frac{\mathcal{R}_N(\mathcal{W}) + \sqrt{\log(1/\delta)}}{\rho \sqrt{N}} \right)
\end{equation}
This establishes the generalization bound of Proposition~\ref{prop:generalization}:
\begin{equation}
\begin{split}
  \mathbb{E}_{x \sim \mathcal{D}} [\mathcal{L}(w; x)] &\le \max_{\|\epsilon\|_2 \le \rho} \frac{1}{N} \sum_{i=1}^N \mathcal{L}(w + \epsilon; x_i) \\
  &+ \mathcal{O}\left( \frac{\mathcal{R}_N(\mathcal{W}) + \sqrt{\log(1/\delta)}}{\rho \sqrt{N}} \right)
\end{split}
\end{equation}
which concludes the proof.

\end{document}
